% !TeX program = pdflatex
% Kompilacja: najpierw pdflatex psa.tex (dwukrotnie), potem pdflatex extra.tex (dwukrotnie);
% odesłania do paragrafów umowy (\ref{art:...}) są pobierane z psa.aux przez pakiet xr-hyper.
\documentclass[12pt,a4paper]{report}
\usepackage[T1]{fontenc}
\usepackage[utf8]{inputenc}
\usepackage{xr-hyper}
\usepackage{basilisk_i18n}
\usepackage{lmodern}
\usepackage{microtype}
\usepackage{setspace}
\usepackage{parskip}
\usepackage{needspace}
\usepackage{tabularx}
\usepackage[normalem]{ulem}
\externaldocument{psa}

\BusinessPlanCompany{Basilisk Systems}
\BusinessPlanWatermarkText{Basilisk Systems}
\BusinessPlanWatermarkColor{black!6}
\BusinessPlanWatermarkFontSize{38}{46}
\BusinessPlanWatermarkAngle{38}
\BusinessPlanAuthorsText{Basilisk Systems --- projekt umowy P.S.A.}
\BusinessPlanConfidentialText{Analiza i opinie --- nie stanowią części umowy}
\BusinessPlanFooterText{Basilisk Systems --- analiza projektu umowy P.S.A. i opinie}
\BusinessPlanVersionText{Wersja 0.9.4-C --- analiza i opinie}

\setstretch{1.18}
\setlist[itemize]{topsep=0.2em,itemsep=0.15em,parsep=0.1em,leftmargin=2.0em}
\setlist[enumerate]{topsep=0.25em,itemsep=0.18em,parsep=0.1em,leftmargin=2.2em}
\setlength{\parindent}{0pt}
\setlength{\parskip}{0.55em}
\setlength{\emergencystretch}{4em}
\renewcommand{\arraystretch}{1.18}

\newcounter{paragraf}
\newcommand{\paragraf}[1]{%
  \refstepcounter{paragraf}%
  \phantomsection%
  \addcontentsline{toc}{section}{\S~\theparagraf\ #1}%
  \Needspace{6\baselineskip}%
  \subsection*{\S~\theparagraf\ #1}%
}
\newenvironment{ContractClause}[1]{%
  \paragraf{#1}%
}{%
  \par
}

\newcommand{\field}[1]{\textbf{[#1]}}
\newcommand{\CompanyFull}{Basilisk Systems prosta spółka akcyjna}
\newcommand{\CompanyShort}{Basilisk Systems P.S.A.}
\newcommand{\KSH}{Kodeks spółek handlowych}
\newcommand{\BoilerplateStrike}[1]{\sout{#1}}

\hypersetup{
  pdftitle={Basilisk Systems P.S.A. --- analiza projektu umowy spółki i opinie kancelarii},
  pdfauthor={Basilisk Systems --- projekt umowy P.S.A.}
}

\begin{document}

\BusinessPlanTitlePage
  {Basilisk Systems}
  {Wersja 0.9.4-C --- analiza i opinie}
  {Analiza i opinie}
  {Analiza projektu umowy P.S.A. i opinie kancelarii}
  {Dokument roboczy Założycieli --- nie stanowi części umowy}
  {0.9.4-C}

\tableofcontents
\newpage

\chapter*{Analiza projektu umowy i opinie}
\addcontentsline{toc}{chapter}{Analiza projektu umowy i opinie}

Dokument roboczy Założycieli towarzyszący projektowi umowy spółki (\texttt{psa.tex}). Zawiera: w Części 1 zapis, jak rozwiązywaliśmy wątpliwości projektowe w toku prac nad umową; w Części 2 nasz audyt paragraf po paragrafie; w Części 3 opinie kancelarii prawnej wraz z naszym stanowiskiem. Nic z tego pliku nie stanowi części umowy ani porady prawnej. Odesłania do paragrafów (§) wskazują numery z pliku umowy; odesłania do dokumentów dodatkowych dotyczą pliku \texttt{extra.tex}.

\clearpage
\section*{Część 1 — Jak rozwiązaliśmy wątpliwości}
\addcontentsline{toc}{section}{Część 1 — Jak rozwiązaliśmy wątpliwości}
\label{annex:rozwiazane-watpliwosci}

\textbf{Status dokumentu.} Załącznik ma charakter informacyjny i dokumentuje sposób rozwiązania najważniejszych wątpliwości prawnych podczas prac nad projektem. Nie stanowi części normatywnej umowy. Zestawienie opiera się na stanie prawnym sprawdzonym na dzień \textbf{15 sierpnia 2026 r.}. Odnośniki prowadzą do oficjalnych tekstów aktów prawnych w ELI oraz do orzeczeń Sądu Najwyższego. Aktualny tekst normatywny każdej omawianej kwestii znajduje się w paragrafie wskazanym przez odesłanie \texttt{\\ref}.

\smallskip
\textbf{Podstawowe źródła:}
\begin{itemize}
  \item \href{https://eli.gov.pl/api/acts/DU/2024/18/text.html}{Kodeks spółek handlowych --- tekst jednolity, Dz.U. 2024 poz. 18}: w szczególności art. 300$^{33}$--300$^{39}$, art. 300$^{45}$, art. 300$^{47}$ § 9, art. 300$^{55}$, art. 300$^{58}$, art. 300$^{79}$ oraz art. 300$^{103}$--300$^{107}$;
  \item \href{https://eli.gov.pl/eli/DU/2026/795/ogl}{Kodeks cywilny --- tekst jednolity, Dz.U. 2026 poz. 795}: w szczególności art. 58, 64, 98--101 i 108;
  \item \href{https://eli.gov.pl/eli/DU/2026/468/ogl}{Kodeks postępowania cywilnego --- tekst jednolity, Dz.U. 2026 poz. 468}: w szczególności art. 1047;
  \item \href{https://eli.gov.pl/eli/DU/2026/236/ogl}{Kodeks rodzinny i opiekuńczy --- tekst jednolity, Dz.U. 2026 poz. 236}: w szczególności art. 31, 46, art. 61$^{7}$ § 2 i art. 61$^{8}$ § 2;
  \item \href{https://eli.gov.pl/eli/DU/2026/176/ogl}{ustawa z 23 stycznia 2026 r. o zmianie KSH i niektórych innych ustaw, Dz.U. 2026 poz. 176} --- uwzględniona przy kontroli aktualności i terminów wejścia w życie.
\end{itemize}

\subsection*{1. Dopuszczalny Nabywca: ustalenie prawne zamiast uznaniowego weta}
\textbf{Aktualny tekst umowy:} §~\ref{art:dopuszczalny-nabywca} i §~\ref{art:zbywanie}.

\textbf{Wątpliwość.} Czy Rada Dyrektorów może według własnej oceny uznać legalnego nabywcę za niedopuszczalnego?

\textbf{Przed:}
\begin{quote}\small
„Każda osoba zamierzająca nabyć prawa dotyczące akcji jest zobowiązana, na uzasadnione żądanie Rady Dyrektorów, przedstawić dokumenty i informacje niezbędne do przeprowadzenia weryfikacji prawnej dopuszczalności nabycia.”
\end{quote}

\textbf{Po:}
\begin{quote}\small
„Weryfikację prawnej dopuszczalności nabycia przeprowadza Rada Dyrektorów na podstawie dokumentów przedstawionych przez nabywcę, publicznie dostępnych urzędowych wykazów oraz, w razie potrzeby, opinii prawnej. Rada Dyrektorów stwierdza wynik weryfikacji w formie uchwały lub protokołu. Uznanie nabywcy za Niedopuszczalnego Nabywcę wymaga wskazania konkretnej podstawy prawnej uniemożliwiającej planowane nabycie.”

„Brak wystarczających informacji do zakończenia weryfikacji nie oznacza sam w sobie uznania osoby za Niedopuszczalnego Nabywcę, lecz wstrzymuje zakończenie transakcji do czasu uzyskania informacji koniecznych do ustalenia jej prawnej dopuszczalności.”
\end{quote}

\textbf{Rozwiązanie.} Status nabywcy jest skutkiem obowiązującego prawa, a Rada dokumentuje wynik. Konstrukcja nie ustanawia ogólnego uznaniowego systemu zgody Spółki na legalnego nabywcę.

\textbf{Zmiana: miękka bramka Kryterium Bezpieczeństwa EU/NATO z wykupem ratalnym (wrzesień 2026).} Poprzednie brzmienie §~\ref{art:zbywanie} ust.~3 powtarzało ustawowy mechanizm art.~300$^{39}$ §~3--6 KSH: Spółka musiała w ciągu miesiąca wskazać nabywcę, który płacił Wartość Godziwą jednorazowo w 30 dni, a w braku środków akcjonariusz zbywał akcje swobodnie --- także poza UE/NATO. Bramka zależała więc od gotówki dostępnej w miesiąc. Obecnie, na podstawie art.~300$^{39}$ §~2 KSH, termin wskazania nabywcy wynosi 3 miesiące, nabywcą może być także Spółka w granicach art.~300$^{47}$ KSH, a cena może zostać rozłożona na raty na zasadach przewidzianych dla spłaty spadkobierców (§~\ref{art:dziedziczenie} ust.~4: do 36 rat, oprocentowanie, zabezpieczenie, wymagalność całości przy opóźnieniu); pierwsza rata jest płatna w 30 dni od ostatecznego ustalenia wartości. Jeżeli Spółka nie wskaże nabywcy albo nabywca nie zapłaci pierwszej raty, akcjonariusz zbywa akcje swobodnie zgodnie z art.~300$^{39}$ §~5 KSH. Równocześnie uszczelniono łańcuch: weryfikacja obejmuje Kryterium i oświadczenie o nabyciu na własny rachunek (§~\ref{art:dopuszczalny-nabywca} ust.~4 i~8), Kryterium nie spełnia osoba działająca na rachunek, w interesie albo na polecenie podmiotu spoza Kryterium (ust.~6), akcjonariusz, który utracił Kryterium po nabyciu, zawiadamia Radę i zbywa akcje w 6 miesięcy (ust.~15), brak wymaganej zgody organu wstrzymuje termin zamiast stanowić odmowę (§~\ref{art:zbywanie} ust.~2 lit.~e), a prawo pierwszeństwa wykonuje się po zgodzie Spółki albo równolegle ze zgłoszeniem (§~\ref{art:pierwszenstwo} ust.~1). Alternatywa twarda (gałąź \texttt{psa\_hard}): odmowa z powodu Kryterium wyłącza §~3--6 art.~300$^{39}$ KSH tak jak lock-up, bez obowiązku wskazania nabywcy. Wybór między wersjami nastąpi po otrzymaniu stanowiska kancelarii.

\textbf{Brzmienie alternatywy twardej} (do porównania; w wariancie twardym §~\ref{art:zbywanie} ust.~4 i §~\ref{art:okres-ograniczenia-zbywania} ust.~1 nie odsyłają już do mechanizmu wskazania nabywcy):
\begin{quote}\small
§~\ref{art:zbywanie} ust.~3: „Na podstawie art.~300$^{39}$ §~2 KSH strony postanawiają, że odmowa zgody z przyczyny określonej w ust.~2 lit.~b nie rodzi po stronie Spółki obowiązku wskazania innego nabywcy, a przepisy art.~300$^{39}$ §~3--6 KSH nie znajdują w tym przypadku zastosowania. Akcjonariusz może wówczas zbyć akcje innemu Dopuszczalnemu Nabywcy spełniającemu Kryterium Bezpieczeństwa EU/NATO z zachowaniem pozostałych postanowień niniejszej umowy; nabycie akcji przez osobę lub podmiot niespełniający Kryterium wymaga zgody Walnego Zgromadzenia zgodnie z §~\ref{art:dopuszczalny-nabywca} ust.~6.”

§~\ref{art:dopuszczalny-nabywca} ust.~5 zdanie drugie: „Niezależnie od tego, brak spełnienia Kryterium Bezpieczeństwa EU/NATO przez planowanego nabywcę stanowi podstawę odmowy zgody Spółki na rozporządzenie akcją na zasadach §~\ref{art:zbywanie} ust.~2 lit.~b i ust.~3, chyba że Walne Zgromadzenie wyrazi zgodę zgodnie z ust.~6.”
\end{quote}

\textbf{Podstawa:} \href{https://eli.gov.pl/api/acts/DU/2024/18/text.html}{KSH --- art. 300$^{39}$ oraz przepisy o rejestrze akcjonariuszy}.

\subsection*{2. Zwrotne zbycie: prywatna sprzedaż, nie umorzenie i nie nabycie akcji własnych}
\textbf{Aktualny tekst umowy:} §~\ref{art:zwrotne-zbycie}.

\textbf{Wątpliwość.} Czy obowiązkowe zbycie na rzecz wskazanego nabywcy za 80\% Wartości Godziwej przy Odejściu Zawinionym może zostać potraktowane jak przymusowe umorzenie albo nabycie akcji własnych?

\textbf{Przed:}
\begin{quote}\small
„Nabywcą akcji w wykonaniu niniejszego mechanizmu nie może być Spółka, chyba że nabycie następuje w odrębnym trybie przewidzianym dla nabywania akcji własnych i po spełnieniu wszystkich właściwych wymogów prawa.”
\end{quote}

\textbf{Po:}
\begin{quote}\small
„Nabywcą akcji w wykonaniu niniejszego mechanizmu nie może być Spółka ani osoba trzecia działająca na rachunek Spółki, chyba że nabycie następuje w odrębnym trybie zgodnym z art. 300$^{47}$ KSH i po spełnieniu wszystkich właściwych wymogów prawa.”
\end{quote}

\textbf{Rozwiązanie.} Akcje nie są unicestwiane, lecz przechodzą na osobę trzecią. Minimalna Wartość Godziwa z art. 300$^{45}$ § 2 KSH dotyczy przymusowego umorzenia. Art. 300$^{47}$ § 9 KSH obejmuje reżimem akcji własnych również osobę trzecią działającą na rachunek Spółki, dlatego została ona wyraźnie wyłączona z mechanizmu prywatnego zbycia.

\textbf{Podstawa:} \href{https://eli.gov.pl/api/acts/DU/2024/18/text.html}{KSH --- art. 300$^{45}$ i art. 300$^{47}$, w szczególności § 9}.

\subsection*{3. Wykonanie zwrotnego zbycia bez dalszego podpisu akcjonariusza}
\textbf{Aktualny tekst umowy:} §~\ref{art:zwrotne-zbycie} i §~\ref{art:breach}.

\textbf{Wątpliwość.} Czy odmowa podpisania dokumentów przez Akcjonariusza Funkcjonalnego unicestwia mechanizm?

\textbf{Przed:}
\begin{quote}\small
„Szczegółowy tryb wykonania obowiązku zbycia [...] określa odrębna umowa wykonawcza [...].”
\end{quote}

\textbf{Po:}
\begin{quote}\small
„W razie niewykonania obowiązku zbycia przez Akcjonariusza Funkcjonalnego Spółka albo uprawniony Dopuszczalny Nabywca może dochodzić wykonania obowiązku zgodnie z prawem, w tym --- jeżeli zachodzą ustawowe przesłanki --- wydania orzeczenia zastępującego wymagane oświadczenie woli zgodnie z art. 64 Kodeksu cywilnego w związku z art. 1047 § 1 Kodeksu postępowania cywilnego. Odrębna umowa wykonawcza powinna określać pełnomocnika w sposób wykluczający niedopuszczalną czynność z samym sobą lub reprezentowanie obu stron bez wymaganej podstawy prawnej.”
\end{quote}

\textbf{Rozwiązanie.} Szczegółowa oferta i ograniczone pełnomocnictwo należą do odrębnej umowy wykonawczej. Art. 101 KC pozwala zrzec się odwołania pełnomocnictwa, gdy uzasadnia to stosunek podstawowy, a art. 108 KC wyznacza granice czynności z samym sobą i reprezentowania obu stron. Gdy wykonanie bezprocesowe nie jest możliwe, art. 64 KC w związku z art. 1047 KPC określa skutek prawomocnego orzeczenia zastępującego wymagane oświadczenie woli, jeżeli istnieje dostatecznie oznaczony obowiązek jego złożenia.

\textbf{Orzecznictwo:}
\begin{itemize}
  \item \href{https://www.sn.pl/sites/orzecznictwo/OrzeczeniaHTML/v\%20csk\%20522-18-1.docx.html}{SN, wyrok z 23 czerwca 2020 r., V CSK 522/18} --- znaczenie stosunku podstawowego dla zrzeczenia się odwołania pełnomocnictwa oraz granic umocowania;
  \item \href{https://www.sn.pl/sites/orzecznictwo/OrzeczeniaHTML/ii\%20cskp\%20593-22.docx.html}{SN, wyrok z 9 grudnia 2022 r., II CSKP 593/22} --- art. 64 KC nie kreuje obowiązku, ale orzeczenie może zastąpić oświadczenie wymagane do wykonania istniejącego obowiązku; art. 1047 KPC reguluje skutek procesowy.
\end{itemize}

\textbf{Podstawa ustawowa:} \href{https://eli.gov.pl/eli/DU/2026/795/ogl}{KC --- art. 64, 98--101 i 108}; \href{https://eli.gov.pl/eli/DU/2026/468/ogl}{KPC --- art. 1047}; \href{https://eli.gov.pl/api/acts/DU/2024/18/text.html}{KSH --- art. 300$^{34}$ § 4 i art. 300$^{36}$--300$^{38}$}.

\subsection*{4. Małżonek: rozdzielenie sfery majątkowej od korporacyjnej}
\textbf{Aktualny tekst umowy:} §~\ref{art:malzonek} i §~\ref{art:dopuszczalny-nabywca}.

\textbf{Wątpliwość.} Czy wejście akcji do majątku wspólnego albo samo ustanie wspólności czyni drugiego małżonka akcjonariuszem i czy potrzebny jest osobny przymusowy wykup za 100\% Wartości Godziwej?

\textbf{Przed:}
\begin{quote}\small
„Jeżeli osoba, o której mowa w ust. 3, jest Niedopuszczalnym Nabywcą, jest zobowiązana [...] zbyć akcje na rzecz Dopuszczalnego Nabywcy [...].”

„Jeżeli osoba [...] nie wykona obowiązku [...], akcje podlegają [...] obowiązkowemu zbyciu [...]. W przypadku obowiązkowego zbycia cena odpowiada 100\% Wartości Godziwej.”
\end{quote}

\textbf{Po:}
\begin{quote}\small
„Sam fakt objęcia akcji wspólnością majątkową małżeńską albo ustania tej wspólności nie powoduje sam przez się, wobec Spółki, nabycia przez drugiego małżonka statusu akcjonariusza ani samodzielnego prawa wykonywania praw korporacyjnych.”

„Jeżeli wskutek prawomocnego orzeczenia sądu, ugody, umowy o podział majątku albo innego prawnie skutecznego zdarzenia małżonek albo były małżonek skutecznie nabędzie akcje, stosuje się do niego na zasadach ogólnych postanowienia niniejszej umowy dotyczące Dopuszczalnego Nabywcy oraz --- jeżeli po skutecznym nabyciu powstanie odpowiednia podstawa prawna --- Prawnie Niedopuszczalnego Posiadacza.”
\end{quote}

\textbf{Rozwiązanie.} Usunięto szczególny mechanizm przymusowego zbycia tylko z powodu statusu małżonka albo byłego małżonka. Prawo rodzinne rozstrzyga przynależność majątkową, natomiast wykonywanie praw korporacyjnych jest odrębną kwestią. Skuteczne nabycie akcji podlega ogólnym zasadom umowy i bezwzględnie obowiązującemu prawu.

\textbf{Orzecznictwo:}
\begin{itemize}
  \item \href{https://www.sn.pl/sprawy/SitePages/Zagadnienia_prawne_SN.aspx?ItemSID=1718-301f4741-66aa-4980-b9fa-873e90506a11\&ListName=Zagadnienia_prawne\&Rok=2022}{SN, uchwała z 6 października 2022 r., III CZP 109/22} --- ustanowienie rozdzielności majątkowej nie powoduje samo przez się uzyskania przez współmałżonka prawa wykonywania praw korporacyjnych z akcji imiennych nabytych przez drugiego małżonka;
  \item \href{https://www.sn.pl/sites/orzecznictwo/OrzeczeniaHTML/iii\%20cskp\%2065-21-1.docx.html}{SN, postanowienie z 28 maja 2021 r., III CSKP 65/21} --- rozdzielenie przynależności praw do majątku wspólnego od statusu korporacyjnego;
  \item \href{https://www.sn.pl/sprawy/SitePages/Zagadnienia_prawne_SN.aspx?ItemSID=786-16544171-be1b-4089-b74b-413997467af2\&ListName=Zagadnienia_prawne\&Rok=2016}{SN, uchwała z 7 lipca 2016 r., III CZP 32/16} --- udział nabyty ze środków wspólnych może należeć do majątku wspólnego, mimo że wspólnikiem jest tylko małżonek dokonujący nabycia.
\end{itemize}

\textbf{Podstawa ustawowa:} \href{https://eli.gov.pl/eli/DU/2026/236/ogl}{KRO --- art. 31 i 46}; \href{https://eli.gov.pl/api/acts/DU/2024/18/text.html}{KSH --- art. 300$^{37}$--300$^{38}$}.

\subsection*{5. Prawnie Niedopuszczalny Posiadacz: zmiana stanu prawnego po legalnym nabyciu}
\textbf{Aktualny tekst umowy:} §~\ref{art:dopuszczalny-nabywca}.

\textbf{Wątpliwość.} Definicja Niedopuszczalnego Nabywcy badała legalność nabycia w chwili transakcji, lecz nie obejmowała sytuacji, gdy akcjonariusz legalnie nabył akcje, a dopiero później powstał wiążący zakaz lub ograniczenie.

\textbf{Przed:}
\begin{quote}\small
„Niedopuszczalnym Nabywcą” jest osoba lub podmiot, którego nabycie akcji jest w chwili planowanego nabycia objęte bezwzględnie obowiązującym zakazem albo ograniczeniem prawnym uniemożliwiającym dokonanie transakcji [...].
\end{quote}

\textbf{Po:}
\begin{quote}\small
„Prawnie Niedopuszczalnym Posiadaczem” jest akcjonariusz, który nabył akcje skutecznie i zgodnie z prawem, lecz po ich nabyciu powstał dotyczący tej osoby albo wykonywania praw z akcji bezwzględnie obowiązujący zakaz lub ograniczenie prawne dotyczące dalszego posiadania akcji, wykonywania praw z akcji, otrzymywania świadczeń z akcji albo rozporządzania akcjami [...].”
\end{quote}

\textbf{Rozwiązanie.} Rozdzielono legalność wejścia do Spółki od późniejszej zmiany stanu prawnego. Status nie jest sankcją korporacyjną ani automatycznym wykupem. Rada dokumentuje konkretną podstawę prawną, a zbycie następuje tylko wtedy, gdy jest prawnie dopuszczalnym sposobem usunięcia stanu niedopuszczalności.

\textbf{Podstawa:} \href{https://eli.gov.pl/api/acts/DU/2024/18/text.html}{KSH --- art. 300$^{33}$--300$^{38}$}; \href{https://eli.gov.pl/eli/DU/2026/795/ogl}{KC --- w szczególności art. 58}.

\subsection*{6. Konflikt interesów: krąg rodzinny, quorum i szczególna reprezentacja}
\textbf{Aktualny tekst umowy:} §~\ref{art:konkurencja} i §~\ref{art:reprezentacja}.

\textbf{Wątpliwość.} Czy próg 2\%, krąg osób bliskich, wyłączenie dyrektora, quorum i reprezentacja przy transakcji z dyrektorem są ze sobą spójne?

\textbf{Przed:}
\begin{quote}\small
„Podmiotem Powiązanym z Dyrektorem” jest podmiot, w którym dyrektor, jego małżonek, osoba bliska albo podmiot przez nich kontrolowany [...].

„Dyrektor pozostający w konflikcie nie uczestniczy w rozstrzyganiu sprawy i jego głosu nie uwzględnia się przy ustalaniu wyniku głosowania. Nie narusza to wymogów quorum wynikających z art. 300$^{58}$ KSH.”
\end{quote}

\textbf{Po:}
\begin{quote}\small
„Na potrzeby niniejszego paragrafu `Osobą Bliską' jest współmałżonek dyrektora oraz jego krewny lub powinowaty do drugiego stopnia włącznie. Stopień pokrewieństwa oraz linię i stopień powinowactwa ustala się zgodnie z art. 61$^{7}$ § 2 oraz art. 61$^{8}$ § 2 Kodeksu rodzinnego i opiekuńczego. Osoba pozostająca z dyrektorem w stosunku pokrewieństwa lub powinowactwa dalszym niż drugi stopień nie jest Osobą Bliską wyłącznie z tego tytułu. Powyższe nie wyłącza zastosowania art. 300$^{55}$ § 1 KSH do osoby, z którą dyrektor jest powiązany osobiście [...].”

„Wyłączenie dyrektora z udziału w rozstrzyganiu sprawy nie zmniejsza pełnego składu Rady Dyrektorów stanowiącego podstawę obliczenia quorum wymaganego przez art. 300$^{58}$ § 2 KSH.”

„Jeżeli wskutek wyłączenia dyrektorów pozostających w konflikcie Rada Dyrektorów nie może podjąć wymaganej uchwały, zawarcie transakcji wymaga uprzedniej uchwały Walnego Zgromadzenia [...]. Uchwała Walnego Zgromadzenia nie zastępuje szczególnego sposobu reprezentacji Spółki wynikającego z art. 300$^{79}$ KSH.”
\end{quote}

\textbf{Rozwiązanie.} Ustawowy test konfliktu pozostaje jakościowy i działa niezależnie od umownego progu 2\%. Automatyczny krąg rodzinny odpowiada art. 300$^{55}$ § 1 KSH, a sposób liczenia stopnia wynika z art. 61$^{7}$ § 2 i art. 61$^{8}$ § 2 KRO. Quorum pozostaje liczone według art. 300$^{58}$ § 2 KSH. Szczególna reprezentacja w umowie lub sporze z dyrektorem podlega art. 300$^{79}$ KSH.

\textbf{Podstawa:} \href{https://eli.gov.pl/api/acts/DU/2024/18/text.html}{KSH --- art. 300$^{55}$ § 1, art. 300$^{58}$ § 2 i art. 300$^{79}$}; \href{https://eli.gov.pl/eli/DU/2026/236/ogl}{KRO --- art. 61$^{7}$ § 2 i art. 61$^{8}$ § 2}.

\subsection*{7. Seria F: późniejsi współzałożyciele oddzieleni od programu motywacyjnego}
\textbf{Aktualny tekst umowy:} §~\ref{art:seria-f}, §~\ref{art:zwrotne-zbycie} i §~\ref{art:program-motywacyjny}.

\textbf{Wątpliwość.} W jednej z wersji projektu zniknęła odrębna ścieżka dla osoby dołączającej po zawiązaniu Spółki jako faktyczny współzałożyciel. Pozostawały tylko akcje serii A i program motywacyjny serii P.

\textbf{Przed:}
\begin{quote}\small
„Wszystkie 1 000 000 akcji emitowanych przy zawiązaniu Spółki stanowią akcje serii A; akcje programu motywacyjnego są emitowane później jako odrębna seria P [...].”
\end{quote}

\textbf{Po:}
\begin{quote}\small
„Akcje dla późniejszych Założycieli Funkcjonalnych mogą być emitowane jako odrębna seria F, a akcje programu motywacyjnego jako odrębna seria P; emisje serii F i P nie zmieniają liczby akcji serii A i są od siebie niezależne.”

„Na podstawie art. 300$^{103}$ KSH niniejsza umowa przewiduje możliwość wyemitowania, w jednej lub kilku emisjach, nie później niż do dnia 31 grudnia 2036 r., maksymalnie 250 000 [...] zwykłych akcji imiennych serii F. Łączna liczba wyemitowanych i istniejących akcji serii F nie może ponadto przekroczyć 20\% wszystkich akcji Spółki istniejących bezpośrednio po danej emisji serii F.”
\end{quote}

\textbf{Rozwiązanie.} Przywrócono osobną serię F dla późniejszych współzałożycieli. Kwalifikacja osoby jako Założyciela Funkcjonalnego Serii F jest oddzielona od właściwej emisji i umowy objęcia. Każda emisja korzysta z podstawy przewidzianej w aktualnej umowie, przy zachowaniu wymaganej uchwały emisyjnej, umowy objęcia i — gdy jest to potrzebne — kwalifikowanej większości dla pozbawienia prawa poboru. Seria F pozostaje niezależna od puli serii P.

\textbf{Podstawa:} \href{https://eli.gov.pl/api/acts/DU/2024/18/text.html}{KSH --- art. 300$^{103}$, art. 300$^{104}$, art. 300$^{105}$, art. 300$^{106}$ § 2 i art. 300$^{107}$ § 3}.



\clearpage
\section*{Część 2 — Audyt paragraf po paragrafie — zgodność i podstawa prawna}
\addcontentsline{toc}{section}{Część 2 — Audyt paragraf po paragrafie — zgodność i podstawa prawna}
\label{annex:audyt-paragraf-po-paragrafie}

\textbf{Status dokumentu.} Niniejszy dokument ma charakter wyłącznie informacyjny i audytowy. Nie zmienia, nie uzupełnia ani nie interpretuje wiążąco postanowień umowy. W szczególności oznaczenie ryzyka w niniejszym dokumencie \textbf{nie powoduje zmiany zakwestionowanego paragrafu}. Celem jest wskazanie, dlaczego przyjętą konstrukcję można uznać za zgodną z prawem albo gdzie pozostaje konkretna wątpliwość wymagająca dalszej weryfikacji.

\textbf{Zakres audytu.} Audyt odnosi się do tekstu umowy w wersji 0.8. Zmiany wprowadzone w późniejszych rewizjach oraz zmiany wariantowe nie były przedmiotem audytu, dlatego przywołane w nim brzmienia, numery ustępów i parametry poszczególnych paragrafów mogą odbiegać od aktualnego tekstu. Zalecane jest ponowne zlecenie audytu zmienionych paragrafów dla każdego wariantu umowy osobno.

\textcolor{red}{\textbf{Uwaga --- do sprawdzenia przez prawnika przed podpisaniem umowy.} Pięć elementów umowy wymaga odrębnej weryfikacji prawnej i do czasu jej zakończenia należy je traktować jako założenia projektowe, a nie ustalony stan prawny: (1)~\textbf{Kryterium Bezpieczeństwa EU/NATO} (§~\ref{art:dopuszczalny-nabywca} ust.~5--6 oraz odesłania w §~\ref{art:zbywanie} ust.~2 lit.~b, ust.~3 i~7, §~\ref{art:dziedziczenie} ust.~1--2, §~\ref{art:malzonek} ust.~3 i §~\ref{art:rada} ust.~3) --- dopuszczalność i skuteczność ograniczenia kręgu nabywców akcji, spadkobierców, małżonków i dyrektorów według obywatelstwa, siedziby i struktury kontroli w świetle KSH, prawa Unii Europejskiej oraz przepisów o kontroli inwestycji, a także zgodność z wymogami programów i zamawiających, na potrzeby których Kryterium wprowadzono; w szczególności dopuszczalność modyfikacji na podstawie art.~300$^{39}$ §~2 KSH mechanizmu wskazania nabywcy przy odmowie zgody z powodu Kryterium (§~\ref{art:zbywanie} ust.~3 --- termin 3 miesięcy, zapłata ratalna, Spółka jako nabywca w granicach art.~300$^{47}$ KSH) i jej skuteczność wobec podmiotu prowadzącego rejestr akcjonariuszy oraz sądu rejestrowego, obowiązki akcjonariusza, który utracił Kryterium po nabyciu (§~\ref{art:dopuszczalny-nabywca} ust.~15), oraz wyłączenie z Kryterium nabycia na cudzy rachunek (ust.~6); przygotowana jest alternatywa twarda (gałąź \texttt{psa\_hard}), w której odmowa z powodu Kryterium wyłącza §~3--6 tego przepisu bez obowiązku wskazania nabywcy --- do oceny, która wersja jest skuteczniejsza i jak obie będą odbierane przez inwestorów VC; (2)~\textbf{klauzule koncesyjne i certyfikacyjne} (§~\ref{art:cel} ust.~4 i §~\ref{art:pkd} ust.~8) --- aktualność i prawidłowość powołanej podstawy prawnej (ustawa z dnia 13 czerwca 2019 r. o wykonywaniu działalności gospodarczej w zakresie wytwarzania i obrotu materiałami wybuchowymi, bronią, amunicją oraz wyrobami i technologią o przeznaczeniu wojskowym lub policyjnym), zakres wyrobów i technologii objętych obowiązkiem koncesyjnym w odniesieniu do bezzałogowych statków powietrznych i ich oprogramowania oraz relacja do przepisów prawa lotniczego i regulacji Unii Europejskiej o bezzałogowych systemach powietrznych; (3)~\textbf{test Kwalifikowanego Inwestora Finansowego} (§~\ref{art:dopuszczalny-nabywca} ust.~14, w projekcie przekreślony) --- Założyciele skłaniają się do usunięcia wyjątku; do oceny, czy jego zachowanie uprościłoby obsługę zbyć wtórnych i wejścia funduszy bez naruszenia wymogów właścicielskich programów finansowania obronnego, oraz dopuszczalność odstąpienia od wymogu dotyczącego beneficjentów rzeczywistych wobec inwestorów funduszu; (4)~\textbf{uchwała ramowa programu motywacyjnego} (§~\ref{art:program-motywacyjny} ust.~3 i~5) --- dopuszczalność powierzenia Radzie Dyrektorów wykonania transz emisji akcji serii P w granicach jednej uchwały Walnego Zgromadzenia oraz objęcia jedną uchwałą o pozbawieniu prawa poboru wszystkich transz na tle art.~300$^{103}$--300$^{107}$ KSH; (5)~\textbf{Przedsięwzięcia Produkcyjne} (§~\ref{art:przedsiewziecia}) --- skutki udostępnienia technologii partnerom na gruncie kontroli eksportu, dopuszczalność uzależnienia udziału partnera od Kryterium Bezpieczeństwa EU/NATO oraz zgodność klauzul licencyjnych (ograniczenie zakresu, prawa do ulepszeń) z prawem konkurencji, w szczególności z przepisami o porozumieniach dotyczących transferu technologii.}

\textbf{Stan prawny audytu: 18 sierpnia 2026 r.} Punktem odniesienia jest w szczególności \href{https://eli.gov.pl/api/acts/DU/2024/18/text.html}{Kodeks spółek handlowych — tekst jednolity Dz.U. 2024 poz. 18 z późniejszymi zmianami}, \href{https://eli.gov.pl/eli/DU/2026/795/ogl}{Kodeks cywilny — Dz.U. 2026 poz. 795}, \href{https://eli.gov.pl/eli/DU/2026/468/ogl}{Kodeks postępowania cywilnego — Dz.U. 2026 poz. 468} oraz \href{https://eli.gov.pl/eli/DU/2026/236/ogl}{Kodeks rodzinny i opiekuńczy — Dz.U. 2026 poz. 236}. Ustawa z 23 stycznia 2026 r. zmieniająca KSH została ogłoszona 17 lutego 2026 r.; jej zasadniczy termin wejścia w życie przypada na 18 lutego 2027 r., dlatego audyt nie stosuje przepisów, które na dzień audytu nie weszły jeszcze w życie. \href{https://eli.gov.pl/eli/DU/2026/176/ogl}{Tekst nowelizacji w ELI}.

\textbf{Metoda.} Ocenie poddano 35 paragrafów normatywnych umowy. Fragmenty widocznie przekreślone w tekście umowy traktowane są jako przeznaczone do usunięcia i nie stanowią podstawy pozytywnej oceny. Oznaczenia mają następujące znaczenie:
\begin{itemize}
  \item \textbf{NIE STWIERDZONO KOLIZJI} — konstrukcja mieści się w swobodzie umowy P.S.A. i nie zidentyfikowano sprzeczności z przepisem bezwzględnie obowiązującym;
  \item \textbf{ZGODNY WARUNKOWO / WYMAGA WYKONANIA} — konstrukcja jest dopuszczalna, ale jej skuteczność zależy od prawidłowego uzupełnienia danych, odrębnej umowy, wpisu, uchwały albo innej czynności;
  \item \textbf{DO WERYFIKACJI / RYZYKO} — zidentyfikowano konkretny punkt, w którym obecne brzmienie może być nieskuteczne, niepełne albo kolidować z ustawowym mechanizmem. Tekstu umowy w tym miejscu nie zmieniono;
  \item \textbf{WERYFIKACJA SPECJALISTYCZNA} — korporacyjna konstrukcja nie budzi zasadniczego zastrzeżenia, lecz materialna poprawność wymaga wiedzy z odrębnej dziedziny regulacyjnej.
\end{itemize}

\subsection*{§~\ref{art:firma} — Firma, forma prawna i czas trwania}
\textbf{Status: NIE STWIERDZONO KOLIZJI.} KSH wymaga określenia firmy i siedziby P.S.A., a firma może być obrana dowolnie pod warunkiem użycia oznaczenia „prosta spółka akcyjna”; ustawa wprost dopuszcza skrót „P.S.A.”. Czas nieoznaczony jest dopuszczalny — czas trwania trzeba wskazać w umowie tylko wtedy, gdy jest oznaczony.

\textbf{Podstawa:} \href{https://eli.gov.pl/api/acts/DU/2024/18/text.html}{art. 300$^{5}$ § 1 pkt 1 i 8 oraz art. 300$^{8}$ KSH}.

\subsection*{§~\ref{art:siedziba} — Siedziba i obszar działania}
\textbf{Status: NIE STWIERDZONO KOLIZJI.} Wskazanie Pruszkowa realizuje ustawowy wymóg określenia siedziby. Postanowienia o prowadzeniu działalności w kraju i za granicą oraz o tworzeniu oddziałów i spółek zależnych nie naruszają ustawowego modelu P.S.A.; pozostają podporządkowane ograniczeniom wynikającym z przepisów szczególnych i pozostałych postanowień umowy.

\textbf{Podstawa:} \href{https://eli.gov.pl/api/acts/DU/2024/18/text.html}{art. 300$^{1}$ § 1 i art. 300$^{5}$ § 1 pkt 1 KSH}.

\subsection*{§~\ref{art:cel} — Cel strategiczny Spółki}
\textbf{Status: NIE STWIERDZONO KOLIZJI.} P.S.A. może zostać utworzona w każdym celu prawnie dopuszczalnym. Opis celu strategicznego ogranicza kierunek gospodarczy Spółki, ale nie zastępuje zezwoleń ani regulacji branżowych. Przekreślony ogólny ustęp o obowiązku przestrzegania zezwoleń nie jest potrzebny dla ważności tej konstrukcji.

\textbf{Podstawa:} \href{https://eli.gov.pl/api/acts/DU/2024/18/text.html}{art. 300$^{1}$ § 1 KSH}.

\subsection*{§~\ref{art:pkd} — Przedmiot działalności i PKD 2025}
\textbf{Status: ZGODNY WARUNKOWO / WYMAGA WERYFIKACJI KLASYFIKACYJNEJ.} Określenie przedmiotu działalności jest elementem obowiązkowym umowy P.S.A. i rozbudowany opis działalności spełnia tę funkcję. Audyt korporacyjny nie potwierdza jednak osobno trafności każdego z 19 kodów PKD względem faktycznie wykonywanej działalności; przed złożeniem wniosku do KRS należy wykonać końcową kontrolę klasyfikacyjną według PKD 2025.

\textbf{Podstawa:} \href{https://eli.gov.pl/api/acts/DU/2024/18/text.html}{art. 300$^{5}$ § 1 pkt 2 KSH}; \href{https://eli.gov.pl/eli/DU/2024/1936/ogl}{rozporządzenie Rady Ministrów z 18 grudnia 2024 r. w sprawie PKD}.

\subsection*{§~\ref{art:akcje} — Akcje i kapitał akcyjny}
\textbf{Status: NIE STWIERDZONO KOLIZJI.} Akcje P.S.A. nie mają wartości nominalnej i nie stanowią części kapitału akcyjnego. Umowa prawidłowo określa liczbę, serię, numery i cenę emisyjną akcji serii A oraz nie wpisuje wysokości kapitału akcyjnego do umowy. Ustawa wymaga, aby kapitał akcyjny wynosił co najmniej 1 zł, ale jego konkretna wysokość jest zmienna poza tekstem umowy.

\textbf{Podstawa:} \href{https://eli.gov.pl/api/acts/DU/2024/18/text.html}{art. 300$^{2}$ § 3, art. 300$^{3}$ § 1--2 i art. 300$^{5}$ § 1 pkt 3 KSH}.

\subsection*{§~\ref{art:cap-table} — Objęcie akcji przez Założycieli}
\textbf{Status: ZGODNY WARUNKOWO / WYMAGA UZUPEŁNIENIA DANYCH.} Ustawa wymaga wskazania akcjonariuszy obejmujących poszczególne akcje. Podział 50/20/10/10/10 i oznaczenie serii A są dopuszczalne, lecz przed zawarciem umowy dane identyfikacyjne, numery przypisanych akcji oraz wkłady w Załączniku nr 1 muszą być kompletne. Pozostawienie pól technicznych niewypełnionych oznaczałoby, że dokument nie jest gotowy do podpisania jako umowa założycielska.

\textbf{Podstawa:} \href{https://eli.gov.pl/api/acts/DU/2024/18/text.html}{art. 300$^{5}$ § 1 pkt 3--4 KSH}.

\subsection*{§~\ref{art:wklady} — Wkłady Założycieli}
\textbf{Status: ZGODNY WARUNKOWO / WYMAGA RZETELNEJ WYCENY I PRZENIESIENIA.} Prawa autorskie, sprzęt i zbywalne prawa patentowe mogą być wkładami do P.S.A. i mogą zasilać kapitał akcyjny. Umowa prawidłowo wyłącza z wkładów serii A pracę i usługi, ponieważ nie mogą one być przeznaczone na kapitał akcyjny. Wartości i tytuły prawne muszą jednak zostać realnie wykazane w Załączniku nr 1 i dokumentach przenoszących; znaczne zawyżenie wkładu przeznaczonego na kapitał powoduje ustawową odpowiedzialność.

\textbf{Podstawa:} \href{https://eli.gov.pl/api/acts/DU/2024/18/text.html}{art. 14 § 1--2, art. 300$^{2}$, art. 300$^{3}$, art. 300$^{4}$, art. 300$^{5}$ § 1 pkt 4--5, art. 300$^{9}$ i art. 300$^{10}$ KSH}.

\subsection*{§~\ref{art:program-motywacyjny} — Program motywacyjny — pula 15\% po pełnym rozwodnieniu}
\textbf{Status: NIE STWIERDZONO KOLIZJI.} Umowa może z góry ustanowić podstawę przyszłych zwykłych emisji, jeśli określa maksymalną liczbę akcji i termin emisji. Pula 176 471 akcji P oraz termin 31 grudnia 2036 r. spełniają tę konstrukcję, a każda konkretna emisja nadal wymaga uchwały i umowy objęcia. Zastrzeżenie większości czterech piątych przy pozbawieniu prawa poboru odpowiada KSH.

\textbf{Podstawa:} \href{https://eli.gov.pl/api/acts/DU/2024/18/text.html}{art. 300$^{103}$--300$^{107}$ KSH}.

\subsection*{§~\ref{art:seria-f} — Założyciele Funkcjonalni Serii F — późniejsi współzałożyciele}
\textbf{Status: NIE STWIERDZONO KOLIZJI.} Rozdzielenie kwalifikacji kandydata od prawnej emisji akcji jest prawidłowe: sama uchwała kwalifikacyjna nie tworzy akcji. Maksymalna liczba 250 000 akcji F, dodatkowy limit 20\%, termin 31 grudnia 2036 r., osobna uchwała emisyjna i umowa objęcia odpowiadają mechanizmowi zwykłej emisji na podstawie istniejących postanowień umowy. Praca i usługi mogą być wkładem na akcje P.S.A., lecz nie mogą być przeznaczone na kapitał akcyjny.

\textbf{Podstawa:} \href{https://eli.gov.pl/api/acts/DU/2024/18/text.html}{art. 300$^{2}$ § 2, art. 14 § 1 oraz art. 300$^{103}$--300$^{107}$ KSH}.

\subsection*{§~\ref{art:zwrotne-zbycie} — Akcjonariusze Funkcjonalni i mechanizm zwrotnego zbycia akcji}
\textbf{Status: ZGODNY WARUNKOWO / WYMAGA UMOWY WYKONAWCZEJ.} Umowa nie konstruuje automatycznego umorzenia ani nabycia akcji własnych przez Spółkę, lecz obowiązek sprzedaży na rzecz wskazanego Dopuszczalnego Nabywcy. Ograniczenia rozporządzania i obowiązki związane z akcjami mogą być ujawniane w rejestrze akcjonariuszy. Skuteczność praktyczna mechanizmu zależy jednak od precyzyjnej umowy wykonawczej: oznaczenia zdarzeń, nabywcy, oferty, ceny, terminów i pełnomocnictwa. Orzeczenie zastępujące oświadczenie woli może służyć wykonaniu dostatecznie skonkretyzowanego obowiązku, ale nie naprawia nieprecyzyjnej podstawy zobowiązania.

\textbf{Podstawa:} \href{https://eli.gov.pl/api/acts/DU/2024/18/text.html}{art. 300$^{33}$ § 1 pkt 10--11, art. 300$^{39}$ i art. 300$^{47}$ KSH}; \href{https://eli.gov.pl/eli/DU/2026/795/ogl}{art. 64, 98--101 i 108 KC}; \href{https://eli.gov.pl/eli/DU/2026/468/ogl}{art. 1047 § 1 KPC}; \href{https://www.sn.pl/sites/orzecznictwo/OrzeczeniaHTML/ii%20cskp%20593-22.docx.html}{SN II CSKP 593/22}.

\subsection*{§~\ref{art:dopuszczalny-nabywca} — Dopuszczalny Nabywca, Niedopuszczalny Nabywca i Prawnie Niedopuszczalny Posiadacz}
\textbf{Status: DO WERYFIKACJI / RYZYKO.} Sam rdzeń konstrukcji — badanie wyłącznie prawnej dopuszczalności i brak uznaniowego weta organu — nie koliduje z KSH. Wątpliwość dotyczy jednak ustępu nakładającego bezpośrednio na \emph{osobę dopiero zamierzającą nabyć akcje} obowiązek przedstawiania dokumentów: przed nabyciem osoba trzecia nie jest jeszcze akcjonariuszem związanym umową Spółki. Wymóg dokumentacyjny może działać jako warunek procedury transakcyjnej, lecz bezpośrednia egzekwowalność tego „obowiązku” wobec obcego nabywcy wymaga weryfikacji.

\textbf{Podstawa problemu:} \href{https://eli.gov.pl/api/acts/DU/2024/18/text.html}{art. 300$^{1}$ § 3 oraz art. 300$^{36}$--300$^{39}$ KSH}.

\subsection*{§~\ref{art:zbywanie} — Zbywanie i obciążanie akcji}
\textbf{Status: DO WERYFIKACJI / RYZYKO.} Ograniczenia zbywania i obciążania akcji są co do zasady dopuszczalne. Ta klauzula powtarza jednak problem z §~\ref{art:dopuszczalny-nabywca}: „nabywca” ma przekazywać informacje jeszcze przed nabyciem, mimo że jako osoba trzecia nie jest jeszcze związany umową. Ponadto skuteczność wszystkich ograniczeń musi być oceniana łącznie z ustawowymi regułami art. 300$^{39}$ KSH i sposobem ujawnienia ograniczeń w rejestrze akcjonariuszy.

\textbf{Podstawa problemu:} \href{https://eli.gov.pl/api/acts/DU/2024/18/text.html}{art. 300$^{33}$, art. 300$^{36}$--300$^{40}$ KSH}.

\subsection*{§~\ref{art:okres-ograniczenia-zbywania} — Okres Ograniczenia Zbywania Akcji Założycieli}
\textbf{Status: DO WERYFIKACJI / RYZYKO — ISTOTNE.} KSH pozwala uzależnić zbycie akcji od zgody Spółki lub ograniczyć je w inny sposób. Jeżeli jednak zbycie jest uzależnione od zgody Spółki, art. 300$^{39}$ § 2 odsyła domyślnie do § 3--6, chyba że umowa wyraźnie stanowi inaczej. Przy odmowie zgody Spółka powinna wtedy wskazać innego nabywcę, a umowa powinna określać termin, cenę lub sposób jej ustalenia i termin zapłaty; w braku tych postanowień ustawa przewiduje możliwość zbycia akcji bez ograniczenia. Obecny 12-miesięczny lock-up wymaga zgody WZ, ale nie wyłącza wprost ustawowego mechanizmu i nie zawiera kompletnego mechanizmu nabywcy zastępczego. Intencja „twardego” lock-upu może więc nie zostać osiągnięta w obecnym brzmieniu.

\textbf{Podstawa problemu:} \href{https://eli.gov.pl/api/acts/DU/2024/18/text.html}{art. 300$^{39}$ § 1--6 KSH}.

\subsection*{§~\ref{art:pierwszenstwo} — Prawo pierwszeństwa nabycia akcji}
\textbf{Status: NIE STWIERDZONO KOLIZJI.} KSH wprost pozwala ustanowić w umowie P.S.A. prawo pierwszeństwa. Umowa może ukształtować procedurę odmiennie od modelu ustawowego; wprowadzenie pierwszego etapu dla Spółki jest możliwe wyłącznie w granicach przepisów o nabywaniu akcji własnych, co obecny paragraf zastrzega. Dopiero po niewykorzystaniu procedury dopuszczona jest sprzedaż osobie trzeciej na niekorzystniejszych dla zbywcy warunkach niełamiących prawa pierwszeństwa.

\textbf{Podstawa:} \href{https://eli.gov.pl/api/acts/DU/2024/18/text.html}{art. 300$^{42}$ i art. 300$^{47}$ KSH}.

\subsection*{§~\ref{art:przylaczenie} — Prawo przyłączenia do zbycia}
\textbf{Status: NIE STWIERDZONO KOLIZJI.} Tag-along nie jest obowiązkowym elementem ustawowym P.S.A., ale mieści się w swobodzie kształtowania praw i ograniczeń związanych ze zbywaniem akcji. Paragraf dostatecznie określa zdarzenie uruchamiające, zakres prawa, zasadę równych warunków i terminy. Skuteczność transakcyjna wymaga uwzględnienia tej klauzuli w dokumentacji sprzedaży.

\textbf{Podstawa:} \href{https://eli.gov.pl/api/acts/DU/2024/18/text.html}{art. 300$^{1}$ § 3, art. 300$^{33}$ § 1 pkt 10--11 oraz art. 300$^{39}$ KSH}.

\subsection*{§~\ref{art:przymusowe-wspolzbycie} — Prawo przymusowego współzbycia}
\textbf{Status: NIE STWIERDZONO KOLIZJI.} Drag-along jest umownym obowiązkiem akcjonariusza związanym z określonym zdarzeniem właścicielskim. Obecna klauzula określa próg 75\%, sprzedaż 100\% akcji, równość ceny dla tego samego rodzaju akcji i ograniczenia odpowiedzialności mniejszości. Nie stwierdzono przepisu KSH zakazującego takiej konstrukcji; praktyczne wykonanie nadal wymaga zachowania formy zbycia akcji i rejestru akcjonariuszy.

\textbf{Podstawa:} \href{https://eli.gov.pl/api/acts/DU/2024/18/text.html}{art. 300$^{1}$ § 3, art. 300$^{33}$ § 1 pkt 10--11, art. 300$^{36}$--300$^{39}$ KSH}.

\subsection*{§~\ref{art:dziedziczenie} — Dziedziczenie akcji}
\textbf{Status: DO WERYFIKACJI / RYZYKO — ISTOTNE.} KSH traktuje dziedziczenie jako przejście akcji z mocy prawa, a jednocześnie pozwala umowie ograniczyć lub wyłączyć wstąpienie spadkobierców do Spółki pod warunkiem określenia zasad ich spłaty. Obecny paragraf nie formułuje wprost ograniczenia lub wyłączenia wstąpienia na podstawie art. 300$^{41}$, lecz najpierw uzależnia wykonywanie praw od weryfikacji Dopuszczalnego Nabywcy, a następnie nakazuje sprzedaż po negatywnej weryfikacji albo po nieprzedstawieniu informacji w 90 dni. Trzeba zweryfikować, czy ten mechanizm dostatecznie wykorzystuje art. 300$^{41}$ i czy 90-dniowy obowiązek sprzedaży jest skuteczny wobec spadkobiercy, który nabył akcje ex lege. Ustęp o wspólnym przedstawicielu kilku spadkobierców jest zgodny z ustawowym modelem współuprawnienia.

\textbf{Podstawa problemu:} \href{https://eli.gov.pl/api/acts/DU/2024/18/text.html}{art. 300$^{37}$ § 2, art. 300$^{38}$ § 1 i 3 oraz art. 300$^{41}$ KSH}.

\subsection*{§~\ref{art:malzonek} — Stosunki majątkowe małżeńskie i akcje}
\textbf{Status: NIE STWIERDZONO KOLIZJI.} Paragraf rozdziela przynależność majątkową akcji od wykonywania praw korporacyjnych i nie próbuje umownie zmieniać skutków wspólności majątkowej. Jeżeli powstaje współuprawnienie, odsyła do ustawowego wspólnego przedstawiciela, a po skutecznym nabyciu przez małżonka stosuje ogólne reguły dotyczące akcji. Takie rozdzielenie odpowiada linii orzeczniczej Sądu Najwyższego dotyczącej różnicy między majątkową przynależnością akcji a wykonywaniem praw korporacyjnych.

\textbf{Podstawa:} \href{https://eli.gov.pl/eli/DU/2026/236/ogl}{KRO, w szczególności art. 31 i 46}; \href{https://eli.gov.pl/api/acts/DU/2024/18/text.html}{art. 300$^{38}$ KSH}; \href{https://www.sn.pl/sprawy/SitePages/Zagadnienia_prawne_SN.aspx?ItemSID=1718-301f4741-66aa-4980-b9fa-873e90506a11&ListName=Zagadnienia_prawne&Rok=2022}{SN III CZP 109/22}.

\subsection*{§~\ref{art:wycena} — Wartość Godziwa i Dzień Wyceny}
\textbf{Status: ZGODNY WARUNKOWO.} Strony mogą umownie zdefiniować standard wyceny używany w kontraktowych mechanizmach zbycia i powierzyć wycenę niezależnemu ekspertowi. Definicja ta nie może jednak zastępować szczególnej ustawowej metody lub minimalnego poziomu spłaty, jeżeli konkretny przypadek podlega bezwzględnie obowiązującemu przepisowi KSH. W obecnej umowie definicja jest używana przede wszystkim dla mechanizmów umownych, co ogranicza to ryzyko.

\textbf{Podstawa:} \href{https://eli.gov.pl/eli/DU/2026/795/ogl}{art. 353$^{1}$ KC}; przykładowe ustawowe odniesienie do wartości godziwej: \href{https://eli.gov.pl/api/acts/DU/2024/18/text.html}{art. 300$^{45}$ § 2 KSH}.

\subsection*{§~\ref{art:organy} — Organy Spółki}
\textbf{Status: NIE STWIERDZONO KOLIZJI.} P.S.A. może działać w systemie monistycznym z Radą Dyrektorów zamiast klasycznego zarządu i rady nadzorczej. Wskazanie Walnego Zgromadzenia oraz Rady Dyrektorów odpowiada ustawowej konstrukcji organów P.S.A.

\textbf{Podstawa:} \href{https://eli.gov.pl/api/acts/DU/2024/18/text.html}{art. 300$^{52}$ oraz art. 300$^{73}$ KSH}.

\subsection*{§~\ref{art:rada} — Rada Dyrektorów}
\textbf{Status: NIE STWIERDZONO KOLIZJI.} Rada Dyrektorów ustawowo łączy prowadzenie spraw, reprezentację i stały nadzór. Umowa może określać liczbę dyrektorów, zasady powoływania, kadencję, podział na dyrektorów wykonawczych i niewykonawczych oraz surowsze reguły wewnętrzne. Delegacja bieżącego prowadzenia spraw nie może obejmować ustawowo zastrzeżonych decyzji strategicznych, a obecny dokument zachowuje to rozdzielenie.

\textbf{Podstawa:} \href{https://eli.gov.pl/api/acts/DU/2024/18/text.html}{art. 300$^{56}$, art. 300$^{58}$ oraz art. 300$^{73}$--300$^{76}$ KSH}.

\subsection*{§~\ref{art:reprezentacja} — Reprezentacja Spółki}
\textbf{Status: NIE STWIERDZONO KOLIZJI.} Zasada dwóch dyrektorów albo dyrektora łącznie z prokurentem jest dopuszczalnym sposobem reprezentacji. Powołanie prokurenta wymaga zgody wszystkich dyrektorów, a odwołać prokurę może każdy dyrektor — obecny tekst odpowiada art. 300$^{75}$ § 3. Dla umów i sporów z dyrektorem zachowano szczególną reprezentację z art. 300$^{79}$, w tym wariant dla dyrektora wykonawczego i spółki jednoosobowej.

\textbf{Podstawa:} \href{https://eli.gov.pl/api/acts/DU/2024/18/text.html}{art. 300$^{75}$ § 3, art. 300$^{78}$ i art. 300$^{79}$ KSH}.

\subsection*{§~\ref{art:rada-zastrzezone} — Sprawy wymagające uchwały Rady Dyrektorów}
\textbf{Status: NIE STWIERDZONO KOLIZJI.} KSH wprost wymaga uchwały Rady przy decyzjach strategicznych, planach biznesowych i ustalaniu struktury organizacyjnej. Umowa może rozszerzyć katalog spraw wymagających uchwały, wprowadzać progi wartościowe oraz dodatkowy wymóg pozytywnego głosu określonego dyrektora, o ile nie narusza ustawowych kompetencji innych organów. Klauzula emisyjna prawidłowo pozostawia samą decyzję o emisji akcjonariuszom.

\textbf{Podstawa:} \href{https://eli.gov.pl/api/acts/DU/2024/18/text.html}{art. 300$^{75}$ § 1--2 i art. 300$^{76}$ KSH}.

\subsection*{§~\ref{art:wz} — Walne Zgromadzenie i sposób obliczania głosów}
\textbf{Status: DO WERYFIKACJI / RYZYKO.} Zdalne uczestnictwo i elektroniczne głosowanie są w P.S.A. dopuszczalne, a umowa może ustanowić surowsze większości. Dwa fragmenty wymagają jednak doprecyzowania prawnego: po pierwsze, art. 300$^{92}$ reguluje \emph{uczestnictwo} elektroniczne w Walnym Zgromadzeniu, podczas gdy art. 300$^{88}$ nadal wskazuje miejsce Walnego Zgromadzenia; sformułowanie, że WZ może „odbywać się ... z wykorzystaniem środków komunikacji elektronicznej”, może być odczytane jako całkowicie wirtualne WZ bez miejsca fizycznego. Po drugie, art. 300$^{87}$ wskazuje ustawowe kanały zwołania, w tym e-mail wpisany do rejestru akcjonariuszy; samo wskazanie adresu jedynie Spółce może nie wystarczyć, jeśli nie spełnia ustawowego sposobu doręczenia. Pozostała część, w tym uchwały poza WZ drogą elektroniczną, ma podstawę w art. 300$^{80}$.

\textbf{Podstawa problemu:} \href{https://eli.gov.pl/api/acts/DU/2024/18/text.html}{art. 300$^{80}$, art. 300$^{87}$--300$^{88}$, art. 300$^{92}$ i art. 300$^{98}$ KSH}.

\subsection*{§~\ref{art:walne-zastrzezone} — Sprawy Zastrzeżone dla Walnego Zgromadzenia}
\textbf{Status: NIE STWIERDZONO KOLIZJI.} Ustawa określa minimalny katalog uchwał akcjonariuszy, a umowa może rozszerzyć go i ustanowić surowsze progi. Większość 75\% wszystkich głosów jest surowsza niż ustawowe reguły dla wielu uchwał i co do zasady dopuszczalna. Paragraf zachowuje szczególne ustawowe wymagania, w tym cztery piąte dla pozbawienia prawa poboru oraz możliwość wymagania zgody akcjonariusza, gdy zmiana zwiększa jego świadczenia lub uszczupla indywidualne prawa.

\textbf{Podstawa:} \href{https://eli.gov.pl/api/acts/DU/2024/18/text.html}{art. 300$^{81}$, art. 300$^{98}$ oraz art. 300$^{106}$ § 2 KSH}.

\subsection*{§~\ref{art:wlasnosc-intelektualna} — Własność intelektualna i aktywa technologiczne}
\textbf{Status: NIE STWIERDZONO KOLIZJI W WARSTWIE KORPORACYJNEJ.} Umowa tworzy obowiązek zapewnienia Spółce praw do kluczowych aktywów, ale prawidłowo przyznaje, że sama ogólna klauzula nie zastępuje dokumentów rozporządzających wymaganych przez prawo autorskie, patentowe lub inne przepisy. To istotne: skuteczność konkretnego przeniesienia zależy od prawidłowej identyfikacji prawa, formy i zakresu odrębnej umowy.

\textbf{Podstawa korporacyjna:} \href{https://eli.gov.pl/api/acts/DU/2024/18/text.html}{art. 300$^{1}$ § 3 oraz art. 300$^{5}$ § 1 pkt 4 KSH}; szczególne przeniesienia wymagają oceny według właściwych ustaw IP.

\subsection*{§~\ref{art:security} — Poufność, bezpieczeństwo informacji i cyberbezpieczeństwo}
\textbf{Status: NIE STWIERDZONO KOLIZJI.} Akcjonariusze P.S.A. mogą zostać zobowiązani w umowie do określonych świadczeń i zachowań; poufność, zgłaszanie incydentów oraz przestrzeganie polityk bezpieczeństwa mieszczą się w tej konstrukcji. Ograniczenie dostępu do Kluczowej Własności Intelektualnej jest wewnętrzną regułą bezpieczeństwa i nie zastępuje ustawowej oceny eksportowej ani sankcyjnej.

\textbf{Podstawa:} \href{https://eli.gov.pl/api/acts/DU/2024/18/text.html}{art. 300$^{1}$ § 3 KSH}; dla dyrektorów dodatkowo art. 300$^{55}$ § 2 KSH.

\subsection*{§~\ref{art:export} — Kontrola eksportu, sankcje i końcowe zastosowanie}
\textbf{Status: WERYFIKACJA SPECJALISTYCZNA.} Jako klauzula wewnętrznej zgodności postanowienie nie koliduje z konstrukcją P.S.A. i może nakładać na akcjonariuszy obowiązki ostrożniejsze niż minimum ustawowe. Audyt prawa spółek nie może jednak potwierdzić, czy konkretne produkty, kod, dane techniczne, usługi lub transakcje Basilisk podlegają określonej klasyfikacji dual-use, wojskowej, sankcyjnej lub licencyjnej. To wymaga odrębnej weryfikacji przez specjalistę od kontroli eksportu i sankcji; samo postanowienie umowy nie rozstrzyga klasyfikacji produktu.

\textbf{Podstawa korporacyjna:} \href{https://eli.gov.pl/api/acts/DU/2024/18/text.html}{art. 300$^{1}$ § 3 KSH}. \textbf{Zakres regulacyjny:} przepisy UE i prawa polskiego właściwe dla danego produktu, technologii, odbiorcy i transakcji.

\subsection*{§~\ref{art:konkurencja} — Zakaz konkurencji, konflikt interesów i transakcje z podmiotami powiązanymi}
\textbf{Status: DO WERYFIKACJI / RYZYKO — WĄSKI, ALE ISTOTNY.} Definicja konfliktu i Osoby Bliskiej odpowiada art. 300$^{55}$ § 1 KSH; prawidłowo zachowano także szczególną reprezentację z art. 300$^{79}$ i ustawowe quorum liczone od pełnego składu. Wątpliwość dotyczy ust. 8--9: gdy wyłączenie skonfliktowanych dyrektorów uniemożliwia podjęcie uchwały Rady, tekst pozwala zastąpić ją uprzednią uchwałą Walnego Zgromadzenia. W odniesieniu do spraw, dla których \emph{ustawa sama} wymaga uchwały Rady Dyrektorów — w szczególności decyzji strategicznych z art. 300$^{75}$ § 2 — nie ma pewności, że uchwała WZ może skutecznie zastąpić ustawowo wymaganą uchwałę Rady. Ten fallback wymaga weryfikacji; pozostała część paragrafu nie została przez to zakwestionowana.

\textbf{Podstawa problemu:} \href{https://eli.gov.pl/api/acts/DU/2024/18/text.html}{art. 300$^{55}$, art. 300$^{58}$ § 2, art. 300$^{75}$ § 2 i art. 300$^{79}$ KSH}; stopień pokrewieństwa i powinowactwa: \href{https://eli.gov.pl/eli/DU/2026/236/ogl}{art. 61$^{7}$ § 2 i art. 61$^{8}$ § 2 KRO}.

\subsection*{§~\ref{art:transakcje-akcjonariusze} — Transakcje z Akcjonariuszami i podmiotami powiązanymi}
\textbf{Status: WYMAGA OPINII PRAWNIKA.} Paragraf ma umożliwić Spółce korzystanie z usług Akcjonariuszy i podmiotów z nimi powiązanych bez tworzenia automatycznego obowiązku przetargu, przy jednoczesnej ochronie Spółki przed transakcjami niegodziwymi lub zawieranymi w konflikcie interesów. Opinii prawnika wymaga w szczególności doprecyzowanie relacji pojęcia „godziwych warunków” do ustawowych ograniczeń świadczeń na rzecz akcjonariuszy, zakresu i skutków konfliktu interesów dyrektora, właściwego trybu reprezentacji Spółki przy umowie z dyrektorem oraz tego, które większe transakcje powinny wymagać dodatkowej zgody i jak należy ustalić progi, aby nie naruszyć ustawowych kompetencji organów. Należy także potwierdzić, że klauzula o braku automatycznego przetargu nie będzie interpretowana jako wyłączenie obowiązków wynikających z przepisów szczególnych, warunków grantów, dotacji, zamówień lub innych źródeł finansowania. Celem opinii nie jest wprowadzenie obowiązku konkurencyjnego wyboru kontrahenta, lecz potwierdzenie, że przyjęty mechanizm jest skuteczny, uczciwy wobec Spółki i pozostałych akcjonariuszy oraz zgodny z bezwzględnie obowiązującymi przepisami.

\textbf{Do sprawdzenia w szczególności:} przepisy KSH dotyczące świadczeń Spółki na rzecz akcjonariusza, konfliktu interesów dyrektora, reprezentacji Spółki w umowie z dyrektorem oraz podziału kompetencji między Radą Dyrektorów i Walnym Zgromadzeniem; dodatkowo przepisy i warunki szczególne mające zastosowanie do konkretnego finansowania lub zamówienia.

\subsection*{§~\ref{art:finansowanie} — Finansowanie Spółki}
\textbf{Status: NIE STWIERDZONO KOLIZJI.} Umowa może zastrzec, że akcjonariusz nie ma obowiązku dalszego finansowania bez odrębnego zobowiązania. Jest to spójne z zasadą, że akcjonariusze P.S.A. są zobowiązani jedynie do świadczeń określonych w umowie Spółki. Dodatkowe instrumenty dające prawa do akcji, kontroli lub Kluczowej Własności Intelektualnej mogą być zastrzeżone dla uchwały akcjonariuszy jako sprawy właścicielskie.

\textbf{Podstawa:} \href{https://eli.gov.pl/api/acts/DU/2024/18/text.html}{art. 300$^{1}$ § 3--4 oraz art. 300$^{81}$ KSH}.

\subsection*{§~\ref{art:zysk} — Zysk, wypłaty i rezerwy}
\textbf{Status: NIE STWIERDZONO KOLIZJI.} Decyzja o przeznaczeniu zysku należy do akcjonariuszy, a ustawa ogranicza maksymalną kwotę wypłat i wprowadza test zdolności do wykonywania zobowiązań przez sześć miesięcy. Umowa może przyjąć surowszy wewnętrzny próg 75\% dla odstąpienia od polityki reinwestowania w pierwszych 36 miesiącach. Obowiązkowy odpis co najmniej 8\% zysku do osiągnięcia ustawowego poziomu kapitału akcyjnego pozostaje niezależnie wiążący.

\textbf{Podstawa:} \href{https://eli.gov.pl/api/acts/DU/2024/18/text.html}{art. 300$^{15}$--300$^{19}$ oraz art. 300$^{98}$ KSH}.

\subsection*{§~\ref{art:impas} — Rozwiązywanie Impasu Decyzyjnego}
\textbf{Status: NIE STWIERDZONO KOLIZJI.} Negocjacja, mediacja i rozstrzygnięcie technicznych kwestii przez eksperta mogą być umownymi mechanizmami porządkowymi. Kluczowe jest prawidłowe zastrzeżenie, że ekspert nie zastępuje uchwały właściwego organu w sprawie strategicznej, personalnej lub właścicielskiej oraz że tryb impasu nie pozwala ominąć kwalifikowanych większości. Tym samym mechanizm nie przejmuje ustawowych kompetencji organów.

\textbf{Podstawa:} \href{https://eli.gov.pl/eli/DU/2026/795/ogl}{art. 353$^{1}$ KC}; kompetencje organów P.S.A.: \href{https://eli.gov.pl/api/acts/DU/2024/18/text.html}{art. 300$^{73}$--300$^{81}$ KSH}.

\subsection*{§~\ref{art:breach} — Naruszenia i wykonanie obowiązków akcjonariusza}
\textbf{Status: ZGODNY WARUNKOWO / WYMAGA DOKUMENTÓW WYKONAWCZYCH.} Roszczenia o zaniechanie, naprawienie szkody i wykonanie zobowiązania mogą wynikać z Kodeksu cywilnego oraz szczególnych postanowień umowy. Paragraf prawidłowo nie udziela blankietowego pełnomocnictwa w samym statucie, lecz przewiduje odrębną umowę wykonawczą ograniczoną do konkretnej procedury. Skuteczność takiego pełnomocnictwa i ofert zależy od ich dokładnego brzmienia i zgodności m.in. z art. 108 KC.

\textbf{Podstawa:} \href{https://eli.gov.pl/eli/DU/2026/795/ogl}{art. 64, 98--101, 108 i przepisy o odpowiedzialności kontraktowej KC}; \href{https://eli.gov.pl/eli/DU/2026/468/ogl}{art. 1047 § 1 KPC}; pomocniczo \href{https://www.sn.pl/sites/orzecznictwo/OrzeczeniaHTML/ii%20cskp%20593-22.docx.html}{SN II CSKP 593/22}.

\subsection*{§~\ref{art:likwidacja} — Rozwiązanie i likwidacja Spółki}
\textbf{Status: NIE STWIERDZONO KOLIZJI.} Umowa może ustanowić surowszy od ustawowego próg dla uchwały o rozwiązaniu Spółki; 75\% wszystkich głosów jest takim zaostrzeniem. W systemie monistycznym ustawowe odwołania do zarządu poza zakresem wyłączonym przez KSH stosuje się odpowiednio do Rady Dyrektorów lub dyrektorów. Zasady podziału majątku likwidacyjnego mogą być określone w umowie z poszanowaniem bezwzględnie obowiązujących zasad ochrony wierzycieli i praw szczególnych.

\textbf{Podstawa:} \href{https://eli.gov.pl/api/acts/DU/2024/18/text.html}{art. 4 § 2$^{1}$, art. 300$^{98}$ § 2 i art. 300$^{120}$ KSH wraz z przepisami o likwidacji stosowanymi odpowiednio}.

\subsection*{§~\ref{art:final} — Postanowienia końcowe}
\textbf{Status: NIE STWIERDZONO KOLIZJI.} Klauzula salwatoryjna nie próbuje automatycznie „uzdrowić” nieważnego postanowienia, lecz zobowiązuje strony do jego zastąpienia. Zmiana umowy pozostaje podporządkowana formie ustawowej, uchwale i rejestracji. Klauzula mediacyjna zachowuje możliwość czynności pilnych i zabezpieczających. Właściwość sądu siedziby jest zgodna z ustawową właściwością dla sporów ze stosunku spółki tam, gdzie ma ona charakter wyłączny.

\textbf{Podstawa:} \href{https://eli.gov.pl/api/acts/DU/2024/18/text.html}{art. 300$^{102}$ KSH}; \href{https://eli.gov.pl/eli/DU/2026/468/ogl}{art. 40 KPC}; ogólna swoboda umów: \href{https://eli.gov.pl/eli/DU/2026/795/ogl}{art. 353$^{1}$ KC}.

\subsection*{Kontrola techniczna Załączników normatywnych nr 1--3}

\textbf{Załącznik nr 1 — DO UZUPEŁNIENIA PRZED PODPISANIEM.} Dane identyfikacyjne Założycieli, dokładne przyporządkowanie akcji, opis każdego wkładu, wartość i podstawa wyceny muszą być kompletne. Wynika to z wymogów art. 300$^{5}$ § 1 pkt 3--5 KSH i z odpowiedzialności za znaczne zawyżenie wkładu przeznaczonego na kapitał akcyjny z art. 300$^{10}$ KSH.

\textbf{Załącznik nr 2 — DO UZUPEŁNIENIA ALBO JEDNOZNACZNEGO OZNACZENIA BRAKU OSÓB.} Jeżeli którykolwiek Założyciel Pierwotny ma być Akcjonariuszem Funkcjonalnym, musi zostać wskazany wraz z datą rozpoczęcia mechanizmu. Jeżeli żaden — załącznik powinien to jednoznacznie stwierdzać przed podpisaniem.

\textbf{Załącznik nr 3 — DO UZUPEŁNIENIA ALBO JEDNOZNACZNEGO OZNACZENIA BRAKU WYJĄTKÓW.} Załącznik ma dokumentować uprzednio ujawnioną działalność konkurencyjną lub inne zaakceptowane wyjątki; brak wyjątków powinien być zapisany wprost zamiast pozostawienia nieoznaczonych pól.

\subsection*{Podsumowanie audytu}

W obecnym brzmieniu audyt nie stwierdził systemowej sprzeczności konstrukcji P.S.A. z KSH. Najważniejsze punkty wymagające dalszej weryfikacji bez zmieniania tekstu to:
\begin{enumerate}
  \item §~\ref{art:dopuszczalny-nabywca} i §~\ref{art:zbywanie} — bezpośrednie obowiązki nakładane na przyszłego nabywcę przed uzyskaniem przez niego statusu akcjonariusza;
  \item §~\ref{art:okres-ograniczenia-zbywania} — skuteczność „twardego” 12-miesięcznego lock-upu wobec mechanizmu art. 300$^{39}$ § 2--6 KSH;
  \item §~\ref{art:dziedziczenie} — relacja 90-dniowej procedury i obowiązkowego zbycia do przejścia akcji z mocy prawa oraz art. 300$^{41}$ KSH;
  \item §~\ref{art:wz} — sformułowanie o elektronicznym Walnym Zgromadzeniu bez jednoznacznego zachowania miejsca WZ oraz adres e-mail niewpisany do rejestru;
  \item §~\ref{art:konkurencja} — możliwość zastąpienia uchwały Rady Dyrektorów uchwałą Walnego Zgromadzenia, gdy uchwały Rady wymaga sama ustawa.
\end{enumerate}

Żaden z powyższych paragrafów nie został w ramach niniejszego audytu zmieniony.

\clearpage
\section*{Część 3 — Opinie kancelarii}
\addcontentsline{toc}{section}{Część 3 — Opinie kancelarii}

Tutaj wklejamy opinie kancelarii w kolejności otrzymania, z oznaczeniem zgodnym z numeracją korespondencji w pliku \texttt{email.md} (R1, R2, \ldots{} --- odpowiedzi na nasze wiadomości Q1, Q2, \ldots). Każdy wpis zawiera: datę, zakres (A--E z zapytania), treść dosłowną albo załączone memorandum, a pod nią stanowisko i decyzje Założycieli co do poszczególnych uwag. Wcześniejszych wpisów nie edytujemy.

\subsection*{R1 --- \field{data} --- \field{zakres}}
\field{treść opinii}

\textbf{Stanowisko Założycieli.} \field{do uzupełnienia po naradzie}

\end{document}
